% Pertemuan 16 (UAS) Praktikum Pemrograman. Dihasilkan oleh tools/build.py dari tools/konten/p16.py.
\documentclass[t]{beamer}
\usetheme{ITERA}
% Mode catatan pembicara: aktif bila \catatan didefinisikan sebelum \input.
\ifdefined\catatan
  \setbeameroption{show only notes}
  \setbeamertemplate{note page}{\vspace*{20pt}\hspace*{4pt}\begin{minipage}{900pt}
    {\ITERAKickerFont\fontsize{11pt}{14pt}\selectfont\color{ITERAGoldText}CATATAN PEMBICARA \ \ |\ \ SLIDE \insertframenumber}\par\vspace{4pt}
    {\ITERADisplay\bfseries\fontsize{22pt}{28pt}\selectfont\color{ITERAStructure}\insertframetitle}\par\vspace{12pt}
    \fontsize{15pt}{23pt}\selectfont\insertnote\end{minipage}}
\fi
\tikzset{fase/.style={draw=none,fill=#1,rounded corners=3pt,minimum width=158pt,minimum height=118pt,text width=146pt,align=center}}

\renewcommand\ITERAmeeting{Pertemuan 16}
\begin{document}
\SlideJudul{IF25-11002 PRAKTIKUM PEMROGRAMAN}{Pertemuan 16\\Ujian Akhir Semester}{Kisi-kisi, format, tata tertib, dan contoh soal untuk materi Pertemuan 9 sampai 14}{Tim Pengajar}{Tahun 2026. Dilaksanakan serentak dengan UAS Algoritma Pemrograman.}
\note{Slide ini dibahas di akhir pertemuan sebelumnya, bukan pada hari ujian.}

\begin{frame}{Format UAS}
\framesubtitle{Durasi 150 menit, 4 soal, 100 poin}
\small
\begin{tabular}{@{}>{\raggedright\arraybackslash}p{66pt}>{\raggedright\arraybackslash}p{448pt}>{\raggedright\arraybackslash}p{105pt}>{\raggedright\arraybackslash}p{76pt}>{\raggedright\arraybackslash}p{143pt}@{}}
\kepala{Soal} & \kepala{Bentuk} & \kepala{Bagian} & \kepala{Poin} & \kepala{CPMK}\\
1 & Menulis program: array atau string dengan fungsi & A & 25 & CPMK0607\\
\barisgenap 2 & Menulis program: struct dengan pengurutan, pencarian, atau rekursi & A & 25 & CPMK0607\\
3 & Tracing fungsi, array, dan rekursi & B & 25 & CPMK0608\\
\barisgenap 4 & Menemukan, menjelaskan, dan memperbaiki kesalahan & B & 25 & CPMK0608\\
\garisdasar
\end{tabular}
\vspace{10pt}

{\small Bagian A dan Bagian B masing-masing 50 poin, sama seperti semua instrumen di mata kuliah ini.}
\note{Tekankan pembagian 50:50 dan bahwa soal tracing tidak perlu komputer.}
\end{frame}

\begin{frame}{Kisi-kisi UAS}
\framesubtitle{Cakupan Pertemuan 9 sampai 14}
\footnotesize
\begin{tabular}{@{}>{\raggedright\arraybackslash}p{189pt}>{\raggedright\arraybackslash}p{104pt}>{\raggedright\arraybackslash}p{454pt}>{\raggedright\arraybackslash}p{104pt}@{}}
\kepala{Topik} & \kepala{Pertemuan} & \kepala{Indikator yang diuji} & \kepala{Bagian}\\
Array satu dimensi & P9 & statistik, maksimum, pergeseran, batas indeks & A, B\\
\barisgenap Array 2D dan string & P10 & pemrosesan baris dan kolom, operasi karakter dan string & A, B\\
Searching & P11 & linear dan binary search, trace low-high-mid & A, B\\
\barisgenap Sorting & P12 & bubble dan selection sort, isi array per putaran & A, B\\
Fungsi & P13 & parameter, nilai kembali, referensi, array sebagai parameter & A, B\\
\barisgenap Rekursi dan struct & P14 & kasus basis, pohon pemanggilan, array of struct & A, B\\
\garisdasar
\end{tabular}
\note{Kisi-kisi ini juga ada di modul persiapan.}
\end{frame}

\begin{frame}{Tata tertib}
\framesubtitle{Berlaku selama ujian}
\begin{columns}[T]
\begin{column}{0.47\textwidth}
\begin{Blok}{Yang boleh}
Satu lembar catatan A4 tulisan tangan sendiri (dua sisi). GDB Online atau g++ di komputer lab untuk soal A.
\end{Blok}
\end{column}
\begin{column}{0.47\textwidth}
\begin{BlokAlert}{Yang tidak boleh}
AI, internet selain GDB Online, catatan elektronik, berdiskusi, atau membuka file selain lembar jawaban.
\end{BlokAlert}
\end{column}
\end{columns}
\vspace{6pt}
\begin{Catatan}
Soal B dikerjakan di kertas tanpa menjalankan program. Pelanggaran tata tertib membuat nilai UAS 0 dan dilaporkan ke program studi.
\end{Catatan}
\note{Bacakan tata tertib di awal ujian.}
\end{frame}

\begin{frame}{Contoh soal untuk berlatih}
\framesubtitle{Folder 02 Hands-on/P16 - UAS - Contoh Soal - Hands-on}
\small
\begin{tabular}{@{}>{\raggedright\arraybackslash}p{47pt}>{\raggedright\arraybackslash}p{312pt}>{\raggedright\arraybackslash}p{388pt}>{\raggedright\arraybackslash}p{104pt}@{}}
\kepala{No} & \kepala{File} & \kepala{Yang diuji} & \kepala{Bagian}\\
1 & \texttt{soal1-nilai.cpp} & array 2D, fungsi, nilai huruf & Bagian A\\
\barisgenap 2 & \texttt{soal2-inventaris.cpp} & struct, bubble sort, binary search, rekursi & Bagian A\\
3 & \texttt{soal3-tracing.cpp} & isi array sesudah fungsi, pohon pemanggilan & Bagian B\\
\barisgenap 4 & \texttt{soal4-jelaskan.cpp} & temukan tiga kesalahan, jelaskan akibat, perbaiki & Bagian B\\
\garisdasar
\end{tabular}
\vspace{10pt}

{\small Keluaran yang benar ada di \texttt{expected/} dan bisa dicek dengan \texttt{bash cek.sh}. Kunci lengkap disimpan dosen.}
\note{Sarankan mahasiswa mengerjakan contoh soal dengan batas waktu seperti ujian sungguhan.}
\end{frame}

\begin{frame}{Rubrik penilaian ujian}
\framesubtitle{Empat level}
\footnotesize
\begin{tabular}{@{}>{\raggedright\arraybackslash}p{156pt}>{\raggedright\arraybackslash}p{195pt}>{\raggedright\arraybackslash}p{185pt}>{\raggedright\arraybackslash}p{166pt}>{\raggedright\arraybackslash}p{136pt}@{}}
\kepala{Aspek} & \kepala{Sangat baik 85--100} & \kepala{Baik 70--84} & \kepala{Cukup 55--69} & \kepala{Kurang <55}\\
A. Program berjalan & kompilasi bersih, keluaran tepat untuk semua uji & keluaran tepat untuk sebagian besar uji & berjalan, keluaran sering salah & tidak terkompilasi\\
\barisgenap A. Struktur kode & tipe, nama variabel, dan alur rapi & satu kekurangan kecil & sulit dibaca & tidak sesuai soal\\
B. Tracing & semua nilai dan keluaran tepat & satu langkah salah & separuh langkah salah & tidak ada atau acak\\
\barisgenap B. Penjelasan & alasan tepat dengan istilah yang benar & tepat tetapi kurang lengkap & sebagian keliru & tidak ada atau disalin\\
\garisdasar
\end{tabular}
\vspace{8pt}

{\small Skor soal = poin soal × skor level rata-rata aspek yang relevan.}
\note{Rubrik yang sama dipakai untuk UTS dan UAS.}
\end{frame}

\begin{frame}{Cara bersiap}
\framesubtitle{Persiapan}
\begin{columns}[T]
\begin{column}{0.31\textwidth}
\begin{Langkah}{1}{Ulangi}
Modul Belajar 9 sampai 14 dan peta materi di slide P15.
\end{Langkah}
\end{column}
\begin{column}{0.31\textwidth}
\begin{Langkah}{2}{Latih}
Kerjakan empat contoh soal P16 dalam 150 menit tanpa bantuan.
\end{Langkah}
\end{column}
\begin{column}{0.31\textwidth}
\begin{Langkah}{3}{Siapkan}
Lembar catatan A4 tulisan tangan dan datang 15 menit lebih awal.
\end{Langkah}
\end{column}
\end{columns}
\note{Ingatkan jadwal dan ruang ujian.}
\end{frame}

\SlidePenutup{Selamat bersiap}{Terima kasih untuk satu semester belajar bersama. Kerjakan dengan tenang.}
\note{Slide penutup.}
\end{document}
